% =====================================================================
%  Experimental section (Methods + Results) for the VizWiz-Captions study.
%  Paste into your manuscript. Requires: booktabs, amsmath, siunitx, hyperref.
%  Define the \cite keys used below in your .bib:
%    vizwiz, blip, lora, peft, bitsandbytes, qlora, cider, bleu, rouge,
%    clipscore, bertscore, sbert, landingai, git, qwen25vl, llava, florence2,
%    easyocr, transformers.
%  All results are on the COMPLETE validation split (7,542 images). CIDEr is
%  corpus-dependent and is only compared within this fixed evaluation set.
% =====================================================================

\section{Experimental Setup}
\label{sec:setup}

\subsection{Dataset and preprocessing}
\label{sec:data}
We study image captioning on \textbf{VizWiz-Captions}~\cite{vizwiz}, a benchmark of
photographs taken by people who are blind, each annotated with five
crowdsourced reference captions. The images are frequently blurry, poorly framed,
rotated, or dominated by textual content (product labels, screens, documents):
\SI{63}{\percent} of the images carry detectable text according to the dataset's
\texttt{text\_detected} flag. This makes VizWiz a stress test for both visual
robustness and the ability to convey on-image text, which is central to the
assistive use case.

We built a loader on the COCO-style annotations that (i) discards pre-canned
captions (\texttt{is\_precanned}) and rejected captions (\texttt{is\_rejected}),
and (ii) keeps only images retaining a minimum number of valid captions.
Table~\ref{tab:data} summarizes the resulting partitions. The \textbf{test} split
ships without public references (the evaluation server was retired), so all
quantitative results are reported on the \textbf{complete validation split}
(\num{7542} images), overall and split by \texttt{text\_detected} into
\emph{text} (\num{5010}) and \emph{no-text} (\num{2532}) subsets.

\begin{table}[t]
\centering
\caption{VizWiz-Captions partitions after cleaning. ``Used'' is the number of
images retained after discarding pre-canned/rejected captions and enforcing the
minimum-caption constraint.}
\label{tab:data}
\begin{tabular}{lrrrr}
\toprule
Split & Images & Raw caps. & Clean caps. & Used \\
\midrule
train & \num{23431} & \num{117155} & \num{100575} & \num{22866} \\
val   & \num{7750}  & \num{38750}  & \num{33145}  & \num{7542}  \\
test  & \num{8000}  & \multicolumn{2}{c}{--- hidden ---} & \num{8000} \\
\bottomrule
\end{tabular}
\end{table}

\subsection{Computing environment}
\label{sec:env}
All experiments were run on a single consumer laptop GPU (NVIDIA RTX~5060 Laptop,
\SI{8}{\giga\byte} VRAM, Blackwell \texttt{sm\_120}). Because the host operating
system blocked the unsigned native CUDA libraries, the full pipeline was executed
inside WSL2 (Ubuntu, Python~3.12) with PyTorch~2.11 (CUDA~12.8) and 4-bit
quantization via \texttt{bitsandbytes}~\cite{bitsandbytes}. The \SI{8}{\giga\byte}
budget is a hard constraint that shapes every design choice (parameter-efficient
adaptation, 4-bit loading, batch size~1 with gradient accumulation for the larger
VLMs).

\subsection{Base model and low-rank adaptation}
\label{sec:lora}
Our primary model is \textbf{BLIP-base}~\cite{blip}
(\texttt{Salesforce/blip-image-captioning-base}, \SI{225}{M} parameters),
fine-tuned in half precision with Low-Rank Adaptation (LoRA)~\cite{lora} through
the PEFT library~\cite{peft}. We consider two adaptation scopes: (i) \emph{text-only},
injecting LoRA into the decoder self-attention projections (\texttt{query},
\texttt{value}); and (ii) \emph{text+vision}, additionally adapting the ViT
attention (\texttt{qkv}). Unless noted, LoRA uses $\alpha{=}2r$ and dropout~$0.05$;
the text-only baseline uses $r{=}16$ (\SI{1.18}{M} trainable parameters,
\SI{0.52}{\percent} of the model) and the scaled text+vision model uses $r{=}32$.
Models are trained with AdamW at learning rate \num{1e-4}, half precision, and
(where memory requires) non-reentrant gradient checkpointing, for up to four
epochs; decoding uses beam search of width~3 with at most \num{30} new tokens.

\subsection{OCR-conditioned captioning}
\label{sec:ocr}
To test whether explicitly reading on-image text helps, we condition the decoder
on an OCR \emph{prefix}. The decoder input is
\begin{equation}
\texttt{[CLS]} \;\oplus\; \mathbf{t}_{\text{ocr}} \;\oplus\;
\mathbf{t}_{\text{cap}} \;\oplus\; \texttt{[SEP]},
\end{equation}
with $\mathbf{t}_{\text{ocr}}$ the OCR text (truncated to \num{32} tokens) and
$\mathbf{t}_{\text{cap}}$ the target caption (truncated to \num{40} tokens). The
loss is masked over \texttt{[CLS]} and the OCR prefix, so the model learns to
produce the caption given the image \emph{and} the transcribed text. At inference
it generates from $\texttt{[CLS]}\oplus\mathbf{t}_{\text{ocr}}$ and the prefix is
stripped by token length.

We compare two OCR sources: \textbf{EasyOCR}~\cite{easyocr}, a free traditional
engine, and \textbf{Landing~AI ADE}~\cite{landingai}, a paid vision agent that
returns both literal transcriptions and scene descriptions. For the latter we
evaluate two cleanings: \emph{text-only} (literal transcription) and \emph{full}
(description~$+$~text). OCR was pre-computed and cached for \emph{all} text-bearing
validation images (ADE non-empty for \num{4957}/\num{5010}, \SI{98.9}{\percent}),
so the prefix is available throughout evaluation rather than on a subset.

\subsection{Distillation for autonomous inference}
\label{sec:distill}
Because the ADE prefix requires a paid call \emph{at inference time}, we tested
whether its benefit can be internalized into a stand-alone model: BLIP is trained
on ADE-derived pseudo-captions and generates without any prefix. We tried two
teachers: ADE's raw descriptive prose, and a \emph{style-matched} teacher (the
Phase-3 prefix model, whose outputs already follow the terse VizWiz style).

\subsection{Alternative vision--language models}
\label{sec:vlms}
We benchmarked four alternative VLMs under the same budget with 4-bit
QLoRA~\cite{qlora} where applicable: \textbf{GIT-large}~\cite{git}
(\texttt{microsoft/git-large-coco}); \textbf{Qwen2.5-VL-3B}~\cite{qwen25vl}
(images $\le$\num{448}\,px, batch~1, grad.\ accumulation~8); and
\textbf{LLaVA-1.5-7B}~\cite{llava} (LoRA on \texttt{q,k,v,o\_proj}, $r{=}16$,
\num{20000} pairs, two epochs). \textbf{Florence-2-large}~\cite{florence2} was
attempted but its custom modeling code was incompatible with our Transformers
release~\cite{transformers} and it was excluded. For the LLM-style VLMs we also
report a deterministic \emph{first-sentence} truncation of the caption
(Sec.~\ref{sec:res:vlms}).

\subsection{Evaluation metrics and protocol}
\label{sec:metrics}
We report \textbf{CIDEr}~\cite{cider}, \textbf{BLEU-1..4}~\cite{bleu} and
\textbf{ROUGE-L}~\cite{rouge} via \texttt{pycocoevalcap} with an in-house
tokenizer (avoiding a Java dependency; METEOR and SPICE are not reported). Because
$n$-gram metrics reward lexical overlap with short references and penalize
paraphrase and verbosity, we complement them with \textbf{CLIPScore}~\cite{clipscore}
(reference-free image--caption cosine, CLIP ViT-B/32), \textbf{BERTScore}~\cite{bertscore}
(RoBERTa-large $F_1$, baseline-rescaled), and \textbf{SBERT} cosine
similarity~\cite{sbert} (\texttt{all-MiniLM-L6-v2}). All metrics are computed on
the complete validation split. We note that CIDEr's TF-IDF weighting is
corpus-specific, so its absolute values are \emph{not} comparable across different
evaluation sets; all comparisons here are within the same fixed \num{7542}-image set.

\subsection{Implementation note: BLIP prefix decoding}
\label{sec:note}
\texttt{BlipForConditionalGeneration.generate} drops the last token of the supplied
\texttt{input\_ids} (it assumes a trailing \texttt{[SEP]}); passing a bare
$\texttt{[CLS]}\oplus\mathbf{t}_{\text{ocr}}$ prefix therefore discarded the final
prefix token and misaligned the length-based prefix removal, deleting the first
generated word. We fix this by appending a filler \texttt{[SEP]}, so the decoder
conditions on the complete prefix and the output offset is correct. All
prefix-based numbers use the corrected decoding.


\section{Results}
\label{sec:results}

\subsection{Fine-tuning and vision-encoder scaling}
\label{sec:res:ft}
Table~\ref{tab:progression} traces the autonomous (prefix-free) model. Text-only
LoRA adds \num{+14.6} CIDEr over zero-shot, and scaling the training set a further
\num{+7.2}. The largest single gain comes from adapting the \emph{vision encoder}:
adding LoRA to the ViT lifts CIDEr by \num{+15.5} (\num{69.10}$\rightarrow$\num{84.63}),
far more than extra epochs or rank ($r{=}64$ did not help). Our best autonomous
model (text+vision, $r{=}32$, four epochs) reaches \textbf{\num{86.93}} CIDEr.

\begin{table}[t]
\centering
\caption{Autonomous fine-tuning progression (validation, $n{=}7542$, beam${=}3$).
Vision-encoder adaptation is the dominant lever.}
\label{tab:progression}
\begin{tabular}{lrrrr}
\toprule
Model & CIDEr & BLEU-4 & BLEU-1 & ROUGE-L \\
\midrule
BLIP-base (zero-shot)                 & 47.27 & 19.54 & 61.29 & 42.75 \\
\;+ LoRA (text), 20k, 2 ep            & 61.91 & 22.08 & 64.47 & 45.93 \\
\;+ LoRA (text), full, 3 ep           & 69.10 & 24.20 & 66.40 & 47.28 \\
\;+ LoRA (text+vis., $r{=}32$), 2 ep  & 84.63 & 27.60 & 69.60 & 49.64 \\
\textbf{\;+ LoRA (text+vis., $r{=}32$), 4 ep} & \textbf{86.93} & \textbf{27.93} & \textbf{70.15} & \textbf{49.83} \\
\bottomrule
\end{tabular}
\end{table}

\subsection{OCR as a decoder prefix}
\label{sec:res:ocr}
Table~\ref{tab:ocr20k} compares the three OCR sources against a no-OCR control on
the \num{20000}-pair subset. \emph{Literal} OCR does not help: both EasyOCR and the
text-only ADE transcription stay at or below the control, including on text-bearing
images. Only the \emph{full} ADE prefix, which includes a scene \emph{description},
surpasses the control (\num{+8.1} overall, \num{+12.3} on text). The benefit is thus
attributable to the multimodal description of a strong external agent---effectively
distillation from a VLM---not to OCR per se.

\begin{table}[t]
\centering
\caption{OCR-prefix captioning, 20k subset (validation CIDEr). Only the full ADE
description beats the no-OCR control.}
\label{tab:ocr20k}
\begin{tabular}{lrrr}
\toprule
Prefix & All & Text & No-text \\
\midrule
Control (no OCR)            & 61.87 & 61.63 & 59.60 \\
EasyOCR                     & 60.67 & 60.02 & 59.26 \\
Landing AI ADE (text-only)  & 61.73 & 61.70 & 59.46 \\
\textbf{Landing AI ADE (full)} & \textbf{69.99} & \textbf{73.96} & \textbf{59.56} \\
\bottomrule
\end{tabular}
\end{table}

\subsection{Distillation}
\label{sec:res:distill}
Internalizing the ADE signal into an autonomous model failed to transfer the gain
(Table~\ref{tab:distill}). Distilling ADE's verbose prose collapsed quality
(\num{61.52}$\rightarrow$\num{35.96}): the model adopted the descriptive style
everywhere---even on no-text images ($59.29\rightarrow43.49$)---and stopped matching
the short human captions. A style-matched teacher removed the collapse
(\num{35.96}$\rightarrow$\num{60.92}) but only reached parity with the control. The
advantage of the strong agent resides in its \emph{fresh per-image inference}, not
in a handful of pseudo-labels.

\begin{table}[t]
\centering
\caption{Distillation toward an autonomous model (validation CIDEr).}
\label{tab:distill}
\begin{tabular}{lrrr}
\toprule
Model & All & Text & No-text \\
\midrule
Control                           & 61.52 & 61.27 & 59.29 \\
ADE prefix (Phase 3)              & 69.99 & 73.96 & 59.56 \\
Distil.\ ADE prose (Phase 4)      & 35.96 & 31.33 & 43.49 \\
Distil.\ style-matched (Phase 4b) & 60.92 & 60.54 & 58.92 \\
\bottomrule
\end{tabular}
\end{table}

\subsection{Combining ADE with vision-encoder scaling}
\label{sec:res:full}
Since ADE is paid at inference regardless, we ask whether it is redundant with
vision-encoder adaptation or complementary. On the 20k subset with the text+vision
model, the ADE description adds \num{+10.7} CIDEr on text-bearing images over a
vision-only control (\num{79.36}$\rightarrow$\num{90.07})---\emph{larger} than the
gain on the text-only model---so the two signals are complementary. Scaled to the
full training set (Table~\ref{tab:full}), the combination reaches
\textbf{\num{94.77}} CIDEr (\textbf{\num{103.38}} on text), \num{+7.84} over the best
autonomous model (\num{+11.63} on text). Noisy EasyOCR, by contrast, stays
\emph{below} the vision-only model ($-1.83$), confirming that \emph{OCR quality},
not OCR presence, is what matters.

\begin{table}[t]
\centering
\caption{Full-scale models (validation CIDEr). The vision$+$ADE combination is best
overall; the vision-only model is the best \emph{autonomous} one.}
\label{tab:full}
\begin{tabular}{llrrr}
\toprule
Model & Inference & All & Text & No-text \\
\midrule
EasyOCR + vision          & free & 85.10 & 89.33 & 73.47 \\
Vision only               & free & 86.93 & 91.75 & 74.16 \\
\textbf{Vision + ADE (full)} & paid & \textbf{94.77} & \textbf{103.38} & 74.48 \\
\bottomrule
\end{tabular}
\end{table}

\subsection{Alternative VLMs and metrics beyond $n$-grams}
\label{sec:res:vlms}
Table~\ref{tab:vlms} reports all models on the complete validation set under both
$n$-gram and embedding-based metrics. Among autonomous backbones, none surpasses
fine-tuned BLIP at comparable cost: GIT-large (\num{84.54}) trails the BLIP vision
model (\num{86.93}), and Qwen2.5-VL-3B (\num{70.54}, $\sim$\num{14}$\times$ slower)
is far behind after suppressing its non-English token leakage.

LLaVA-1.5-7B is instructive for \emph{evaluation}. Raw, it is the best on-image
reader but pathologically verbose (\num{23.7} words/caption vs.\ $\sim$\num{12} for
humans), degenerating into repetition after an accurate first sentence; this
collapses every metric (CIDEr~\num{15.20}, BERTScore~\num{31.21}). A trivial,
deterministic \emph{first-sentence truncation} turns it into a top autonomous model:
CIDEr~\num{90.15} (above the \num{86.93} vision-only model) and the best score of all
on BERTScore (\num{56.10}) and SBERT (\num{69.76}). The corpus-level metrics never
detected the defect---only inspecting the captions did---underscoring that
\emph{capability and presentation are distinct axes} that aggregate metrics conflate.
CLIPScore is flat ($\sim$\num{0.74}) and fails to discriminate, consistent with CLIP's
poor on-image text reading in this domain. The $n$-gram and semantic metrics also
diverge: \textbf{BLIP+ADE} tops CIDEr but sits near the bottom of the quality models
on BERTScore, whereas the fluent LLM-style captioners (GIT, LLaVA-1st) lead the
semantic metrics---each metric rewards a different notion of quality.

\begin{table*}[t]
\centering
\caption{Model comparison on the complete validation set ($n{=}7542$). BLEU-4,
ROUGE-L, BERTScore and SBERT are $\times 100$; CLIPScore is raw cosine. BLIP/GIT
decode with beam${=}3$; Qwen/LLaVA decode greedily. Best per column in \textbf{bold}.}
\label{tab:vlms}
\begin{tabular}{lcrrrrrr}
\toprule
Model & Inference & CIDEr & BLEU-4 & ROUGE-L & CLIPScore & BERTScore & SBERT \\
\midrule
BLIP + vision                 & free & 86.93 & 27.93 & 49.83 & 0.74 & 50.19 & 67.87 \\
\textbf{BLIP + ADE (full)}    & paid & \textbf{94.77} & \textbf{29.46} & \textbf{50.51} & 0.75 & 50.07 & 69.33 \\
GIT-large                     & free & 84.54 & 28.12 & 49.75 & 0.74 & 52.00 & 69.43 \\
Qwen2.5-VL-3B                 & free & 70.54 & 20.46 & 44.89 & 0.74 & 51.14 & 65.90 \\
LLaVA-1.5-7B (raw)            & free & 15.20 & 12.62 & 33.93 & 0.76 & 31.21 & 60.89 \\
LLaVA-1.5-7B (1st sentence)   & free & 90.15 & 27.77 & 50.04 & 0.75 & \textbf{56.10} & \textbf{69.76} \\
\bottomrule
\end{tabular}
\end{table*}
